\documentclass[11pt,a4paper]{article}
\usepackage[T1]{fontenc}
\usepackage[utf8]{inputenc}
\usepackage{lmodern}
\usepackage[margin=26mm]{geometry}
\usepackage{amsmath,amssymb,amsthm,mathtools,booktabs,array}
\usepackage[hidelinks]{hyperref}
\usepackage{microtype}
\usepackage{float}
\newtheorem{theorem}{Theorem}
\newtheorem{lemma}[theorem]{Lemma}
\newtheorem{proposition}[theorem]{Proposition}
\newtheorem{corollary}[theorem]{Corollary}
\newcommand{\ZZ}{\mathbb Z}
\newcommand{\QQ}{\mathbb Q}
\newcommand{\NN}{\mathbb Z_{\geq0}}
\newcommand{\supp}{\operatorname{supp}}
\newcommand{\Fib}{\mathcal T}
\title{A degree-18 counterexample to the group-order bound\newline for phylogenetic complexity}
\author{Grisha Pochuev}
\date{8 September 2026\\[3pt]\small Working manuscript for independent verification}
\begin{document}
\maketitle
\begin{abstract}
We exhibit two compatible $18\times3$ tables of zero-sum rows over
$G=\ZZ/17\ZZ$. Their column multisets admit exactly these two tables,
up to row permutation, and the tables have disjoint row supports.
Consequently, no sequence of moves involving at most 17 rows connects
them. Thus $\phi(\ZZ/17\ZZ,3)\geq18$, contradicting the proposed bound
$\phi(G)\leq |G|$. The proof consists of explicit data and elementary
rational linear algebra. We also give a direct ideal-theoretic proof,
a boundary-point obstruction even to set-theoretic generation in degrees
at most 17, and reproducible exact checks.
\end{abstract}

\section{Statement and a general obstruction}
A flow of length $n$ over a finite abelian group $G$ is a tuple
$(g_1,\ldots,g_n)$ with $\sum_j g_j=0$ in $G$.
Two tables of flows are \emph{compatible} if their corresponding columns
are the same multisets. A move of size at most $d$ replaces at most $d$
rows with a compatible collection of flows. Tables are considered up to
permutation of their rows.

This is the table formulation of ideal generation in
\cite[Section 2]{MV}. The proposed group-order bound and its three-leaf
special case appear as Conjectures 4.1 and 4.2 there; see also
\cite[Conjecture 27 in the arXiv version]{SS}.

\begin{theorem}\label{thm:main}
The tables $A,B$ in Table~\ref{tab:data} are compatible tables over
$\ZZ/17\ZZ$, each with 18 rows and three columns. Their common column
multisets admit exactly two tables, namely $A$ and $B$, up to row
permutation. Every move changing either table has size 18. In particular,
\[
 \phi(\ZZ/17\ZZ,n)\geq18\quad\text{for every }n\geq3,
 \qquad \phi(\ZZ/17\ZZ)>17.
\]
\end{theorem}

\begin{lemma}\label{lem:two}
Suppose a collection of fixed column multisets admits exactly two
$m$-row tables of flows $A,B$. If $A$ and $B$ have no row type in common,
they cannot be connected by moves of size less than $m$.
\end{lemma}
\begin{proof}
A move of size less than $m$ preserves at least one original row, so
cannot take $A$ to $B$. Every move preserves the column multisets; hence
its result is again $A$ or $B$. No first nontrivial move from $A$ is
possible, which excludes a chain of any length.
\end{proof}

\clearpage
\section{The two explicit tables}
All entries below are residues modulo 17. A row of type $r_i$ is used
$\alpha_i$ times in $A$ and $\beta_i$ times in $B$.
\begin{table}[H]
\centering
\renewcommand{\arraystretch}{1.06}
\begin{tabular}{rcrr}
\toprule
$i$ & $r_i$ & $\alpha_i$ & $\beta_i$\\
\midrule
1 & $(1,3,13)$ & 2 & 0\\
2 & $(1,11,5)$ & 1 & 0\\
3 & $(1,15,1)$ & 0 & 3\\
4 & $(7,3,7)$ & 0 & 2\\
5 & $(7,8,2)$ & 3 & 0\\
6 & $(7,16,11)$ & 0 & 1\\
7 & $(8,8,1)$ & 0 & 2\\
8 & $(8,15,11)$ & 2 & 0\\
9 & $(10,0,7)$ & 1 & 0\\
10 & $(10,6,1)$ & 1 & 0\\
11 & $(10,10,14)$ & 0 & 1\\
12 & $(10,11,13)$ & 0 & 1\\
13 & $(11,1,5)$ & 0 & 2\\
14 & $(11,10,13)$ & 1 & 0\\
15 & $(11,16,7)$ & 1 & 0\\
16 & $(14,1,2)$ & 0 & 2\\
17 & $(14,6,14)$ & 1 & 0\\
18 & $(14,15,5)$ & 1 & 0\\
19 & $(15,0,2)$ & 0 & 1\\
20 & $(15,1,1)$ & 4 & 0\\
21 & $(15,6,13)$ & 0 & 2\\
22 & $(15,8,11)$ & 0 & 1\\
\midrule
 & Total & 18 & 18\\
\bottomrule
\end{tabular}
\caption{A compatible pair with disjoint row supports.}\label{tab:data}
\end{table}
Each row sum is 17 or 34. Write $v^{[m]}$ for $m$ occurrences of $v$.
The common column multisets are
\begin{align*}
 C_1&=\{1^{[3]},7^{[3]},8^{[2]},10^{[2]},11^{[2]},14^{[2]},15^{[4]}\},\\
 C_2&=\{0^{[1]},1^{[4]},3^{[2]},6^{[2]},8^{[3]},10^{[1]},11^{[1]},15^{[3]},16^{[1]}\},\\
 C_3&=\{1^{[5]},2^{[3]},5^{[2]},7^{[2]},11^{[2]},13^{[3]},14^{[1]}\}.
\end{align*}
Let $S_j=\supp C_j$. Every compatible row lies in $S_1\times S_2\times S_3$.
For each $a\in S_1$ and $b\in S_2$, its third entry must be $-a-b$ modulo 17.
Checking membership in $S_3$ gives exactly the 22 types in the table.
For first entries $1,7,8,10,11,14,15$, the numbers of permitted types are,
respectively, $3,3,2,4,3,3,4$.

\clearpage
\section{An elementary kernel computation}
Let $M$ be the $23\times22$ incidence matrix whose rows are labelled by
$(j,v)$, with $j=1,2,3$ and $v\in S_j$ in increasing order, and whose
columns are the types $r_i$:
\[
 M_{(j,v),i}=\begin{cases}1,&(r_i)_j=v,\\0,&\text{otherwise.}\end{cases}
\]
Thus $Mx$ records the column multiplicities of a table with type counts $x$.
All equalities involving $M$ below are over $\QQ$, \emph{not} modulo 17.

\begin{lemma}\label{lem:kernel}
$\ker_{\QQ}M=\QQ(\beta-\alpha)$.
\end{lemma}
\begin{proof}
Let $z\in\ker_{\QQ}M$ and put
\[
 a=z_1,\quad b=z_2,\quad d=z_6,\quad e=z_7,\quad
 f=z_9,\quad g=z_{11},\quad h=z_{17}.
\]
The marginal equations for the first two columns give, by successive
substitution,
\begin{align*}
z=(&a,b,-a-b,-a,a-d,d,e,-e,f,b-f-g,g,-b,\\
   &g+d,-g,-d,-h-a-b-e,h,a+b+e,-f,\\
   &-g-d+h+a+b+e,-b+f+g-h,-a+d-e).
\end{align*}
Conversely, this expression satisfies all first- and second-column
equations for arbitrary values of the seven parameters.
The third-column equations for the values $14,7,11,2,5,13$, in that order,
are
\begin{align*}
 g+h&=0,& -a-d+f&=0,& -a+2d-2e&=0,\\
 -b-d-e-f-h&=0,& a+2b+d+e+g&=0,& a-2b+f-h&=0.
\end{align*}
The first three give $h=-g$, $f=a+d$, and $a=2d-2e$.
The fourth then gives $g=b+4d-e$, while the last two give
\[
 7d-2e+3b=0,\qquad b=9d-5e.
\]
Consequently,
\[
 0=7d-2e+3(9d-5e)=34d-17e,
\]
so $e=2d$. Substitution gives
\[
 (a,b,d,e,f,g,h)=d(-2,-1,1,2,-1,1,-1),
\]
and hence
\begin{align*}
z=d(&-2,-1,3,2,-3,1,2,-2,-1,-1,1,1,\\
    &2,-1,-1,2,-1,-1,1,-4,2,1)=d(\beta-\alpha).
\end{align*}
The remaining third-column equation, for value 1, is
$b-d+2e-f-2g+h=0$ and is satisfied. Since the compatible vectors
$\alpha,\beta$ are distinct, $\beta-\alpha$ is a nonzero element of the
kernel. This proves the assertion.
\end{proof}

\begin{proof}[Proof of Theorem~\ref{thm:main}]
Any compatible table has a count vector $x\in\NN^{22}$ with
$Mx=M\alpha$. By Lemma~\ref{lem:kernel},
\[
 x=\alpha+t(\beta-\alpha),\qquad t\in\QQ.
\]
The sixth coordinate gives $x_6=t\in\NN$, and the second gives
$x_2=1-t\geq0$. Thus $t\in\{0,1\}$ and $x$ is $\alpha$ or $\beta$.
Lemma~\ref{lem:two} excludes every move of size at most 17. Replacing all
18 rows does work. Appending zero columns preserves exactly the same
compatibility problem, proving the assertion for every $n\geq3$.
\end{proof}

\clearpage
\section{A direct ideal-theoretic certificate}
This section independently explains why the obstruction concerns the
full group-based ideal, not merely a chosen subset of moves.
Let $k$ be a field and put
\[
 \mathcal F=\{(a,b,c)\in(\ZZ/17\ZZ)^3:a+b+c=0\},\qquad
 R=k[X_r:r\in\mathcal F].
\]
The three-leaf group-based toric ideal is the kernel $I$ of the map
\[
 X_{(a,b,c)}\longmapsto U_aV_bW_c.
\]
Let $E=\{r_1,\ldots,r_{22}\}$ be the allowed types from Table~\ref{tab:data},
write $X_i=X_{r_i}$, and set $R_E=k[X_1,\ldots,X_{22}]$.
Consider
\begin{align*}
 F=X^\alpha-X^\beta
   ={}&X_1^2X_2X_5^3X_8^2X_9X_{10}X_{14}X_{15}X_{17}X_{18}X_{20}^4\\
     &-X_3^3X_4^2X_6X_7^2X_{11}X_{12}X_{13}^2X_{16}^2X_{19}X_{21}^2X_{22}.
\end{align*}
Both monomials have degree 18, and $F\in I$ by compatibility.

\begin{proposition}\label{prop:retract}
For the specialization $\pi:R\longrightarrow R_E$ which sets all
variables outside $E$ to zero, one has
\[
 \pi(I)=I\cap R_E=(F).
\]
In particular, every element of $I$ of total degree at most 17 specializes
to zero, whereas $\pi(F)=F\ne0$.
\end{proposition}
\begin{proof}
First consider a binomial $X^u-X^v\in I$. If one monomial uses only
variables in $E$, it has no column label outside the sets $S_j$.
Its equal column multiplicities force the other monomial to have this
same property. As $E$ contains \emph{all} flows with labels in $S_j$,
the other monomial also uses only $E$. Thus $\pi$ either retains both
monomials or kills both. The kernel of a monomial map is spanned by
binomials with equal images: this follows simply by grouping terms with
the same image monomial. It follows that $\pi(I)=I\cap R_E$.

For any binomial in this restricted ideal, its exponent difference lies
in $\ker_{\ZZ}M$. Lemma~\ref{lem:kernel}, and the coordinate
$(\beta-\alpha)_6=1$, give
\[
 \ker_{\ZZ}M=\ZZ(\beta-\alpha).
\]
After interchanging the two monomials if necessary and factoring their
greatest common monomial divisor, the binomial is therefore a monomial
multiple of
\[
 (X^\alpha)^m-(X^\beta)^m,
 \qquad m\geq1.
\]
This polynomial is divisible by $F$. Conversely, $F$ belongs to the
restricted ideal. Hence the restricted ideal is $(F)$.
A nonzero polynomial multiple of the homogeneous degree-18 polynomial
$F$ cannot have total degree less than 18, which proves the last claim.
\end{proof}

Were $I$ generated by polynomials of degrees at most 17, their
specializations would generate $\pi(I)$. All would be zero, although
$\pi(I)=(F)\ne0$. This is a direct contradiction. It establishes the
same degree lower bound without invoking a connectivity theorem.

\clearpage
\section{A stronger boundary obstruction and exact checks}
\begin{corollary}
Let $J_{\leq17}\subset I$ be the ideal generated by all elements of $I$
of total degree at most 17. Then
\[
 F\notin\sqrt{J_{\leq17}}.
\]
Thus these low-degree equations do not even define the full variety
set-theoretically.
\end{corollary}
\begin{proof}
Define a point $P$ by $X_r(P)=1$ for the 11 row types occurring in $A$,
and $X_r(P)=0$ for all other flows. This point has all variables outside
$E$ zero. Proposition~\ref{prop:retract} shows that every element of
$J_{\leq17}$ vanishes at $P$, but $F(P)=1-0=1$. No positive power of
$F$ can therefore lie in $J_{\leq17}$.
\end{proof}
This is a boundary point, with many coordinates zero. No claim is made
about the analogous question after restricting to the locus where all
coordinates are nonzero. Also, no equality $\phi(\ZZ/17\ZZ)=18$ is asserted;
the lower bound is sufficient to disprove the proposed inequality.

\subsection*{Reproducible verification}
The accompanying file \texttt{verify\_phylogenetic\_Z17.py} uses only the
Python standard library. It checks all $17^3$ triples to construct the
allowed row types, the row sums and column multiplicities, an exact
rational row reduction, and a full enumeration of compatible count
vectors. It returns exactly the two vectors $\alpha,\beta$.

An additional exact determinant certificate is available. In the matrix
$M$, delete rows 16 and 23 and column 22, numbering from 1 with the
ordering specified above. The resulting $21\times21$ determinant is 17.
Together with $M(\beta-\alpha)=0$, this certifies rank 21. The elementary
proof does not depend on trusting this numerical determinant claim.

The separately written \texttt{audit\_Z17\_independent.py} uses a different
check. It fixes the 18 first-column entries in increasing order and counts
all second- and third-column fillings by dynamic programming on remaining
multiplicities. There are 96 labelled fillings in total. Table $A$
accounts for 24 and table $B$ for 72, by multinomial counting within equal
first-entry blocks. Every compatible unordered table has at least one
such filling, so no third table exists. This check uses neither matrix
rank nor the one-dimensional-kernel conclusion; it visits 981 memoized
states. All counts are exact integers.

\subsection*{Verification status}
The displayed argument is a complete mathematical proof, with finite
arithmetic certificates and independently implemented checks. The note
was prepared with AI assistance. Independent expert review and priority
assessment have not yet been obtained. The material has not been
submitted or published by the assistant.

\begin{thebibliography}{9}
\bibitem{MV}
Mateusz Micha\l ek and Emanuele Ventura,
\emph{Finite phylogenetic complexity and combinatorics of tables},
Algebra \& Number Theory \textbf{11} (2017), 235--252.
\href{https://arxiv.org/abs/1606.07263}{arXiv:1606.07263}.
\bibitem{SS}
Bernd Sturmfels and Seth Sullivant,
\emph{Toric ideals of phylogenetic invariants},
Journal of Computational Biology \textbf{12} (2005), 204--228.
\href{https://arxiv.org/abs/q-bio/0402015}{arXiv:q-bio/0402015}.
Conjecture numbering in this note refers explicitly to the arXiv version.
\end{thebibliography}
\end{document}
